% Folder 33, batch 07 — archivists' pages 121–140.
% Pass: Opus 5.5 (claude-opus-5-5), 2026-10-03 — first pass, unchecked against the pages by a human.
%
% Even pages 122–140 are versos carrying typescripts unrelated to the notes
% (an upside-down French typescript on relative divisors, then pages of an
% English typescript on regular local domains and quadratic transforms) and are
% not transcribed. The hand on the odd pages is fast; prose is heavily marked.
\documentclass[11pt,a4paper]{article}
\input{../preamble/grothendieck}

\folder{33}
\batch{7}
\pages{121}{140}
\foldertitle{SGA 1965. MS [Manuscrit] brouillon : notes manuscrites (s.d.).}
\dating{1965-[vers 1971]}
\watermark{Édition de démonstration}

\begin{document}

\page{121}
\note{la page s'ouvre sur un point numéroté (3) : l'énumération commence avant ce lot.}

(3) Nous en savons assez (formules de Künneth, dualité globale) pour faire la
Lefschetz « \uncertain{cuite et pressée} », en redescendant sur terre. Il semble
que pour les courbes, on \uncertain{récupère} par des arguments élémentaires
(sans utiliser \uncertain{les singularités}) des Lefschetz plus \ill{}, par
faisceaux d'\uncertain{irrégularités} $\rightsquigarrow$ \uncertain{rationalité}
des fonctions $L$.

Nous pouvons aussi, au lieu de cela, continuer en \uncertain{généralisant} les
\uncertain{foncteurs}, \uncertain{en l'occurrence} : \uncertain{formule de
dualité} (= \uncertain{dualité locale}), Th. de bidualité, Lefschetz-\ill{},
\struck{\ill{}} \add{et} \struck{\uncertain{utilisant les fonctions} $L$}
\ill{} est \uncertain{nouveau pavé}.

Étant parti dans un \uncertain{formulaire} de dualité \uncertain{général}, je
\uncertain{préfère} en compléter les lignes essentielles avant de
\struck{\uncertain{me fig}} me mettre aux applications de détail.

\page{123}
\section{(30) Généralités sur Lefschetz. \struck{Lefschetz} Lefschetz-Verdier}

\textbf{Préliminaires à la formule de \struck{\ill{}} Lefschetz.}

Classiquement, on a une variété compacte $X$, un endomorphisme $f$,
\uncertain{d'où} (\ill{} \ill{} \ill{} \ill{} $k$) \uncertain{des
endomorphismes} des $f^{(i)}$ des $H^i(X)$, et la formule dit
\[
  \sum (-1)^i \operatorname{Tr} f^{(i)} = \text{nb des pts fixes de } f
\]
avec grains de sel pour la dernière notion (si les pts fixes \uncertain{ne
sont pas isolés}, il faut une \uncertain{multiplicité} \ill{} \ill{}).
\marginal{\uncertain{Walters}}

En fait, les \uncertain{généralisations} \add{\uncertain{et passages à la
formule}} \uncertain{interprétant} le \ill{} \ill{} \ill{} \ill{} $X$ \ill{}
\ill{} \ill{} \ill{} \ill{} \ill{}), et Verdier a \ill{} \add{\uncertain{par
une méthode et les} \ill{} \ill{} \ill{}} \ill{} \ill{} \ill{} \ill{} un
faisceau \uncertain{généralisé} \ill{} \ill{} \ill{} \ill{} \uncertain{des}
\ill{} coefficients $\pm$ quelconques (\ill{} \ill{} \ill{}
\uncertain{interprétés}). Les \ill{} \ill{} très simples dans des cas
\ill{} \ill{} \ill{} \ill{} \ill{} \ill{} \ill{} \ill{} — maintenant un
\ill{} \ill{} « élémentaire » (variété, \uncertain{coefficients}
\uncertain{localement} constants), \ill{} résultant de la
\uncertain{dualité} de Poincaré \ill{} \ill{} \ill{} \ill{} \ill{} \ill{}
de la \ill{} \ill{} \ill{} \ill{} \ill{}. Dans le cadre des
\uncertain{schémas} \add{(utilisant des \uncertain{outils un peu plus
profonds})} \ill{} Verdier \ill{} sans autres \ill{}, et nous donnerons
\struck{nous} \uncertain{donnerons} \ill{} en détail \ill{} le
« cas élémentaire », qui nous permettra ensuite

\page{125}
des \uncertain{limites} de cas élémentaires de coefficients non constants
\uncertain{algébriques}, en \uncertain{demandant} \uncertain{seulement}
explicitement les \uncertain{faisceaux} locaux.

Une difficulté \add{technique} préliminaire dans la \ill{} \ill{} \ill{} de
\uncertain{tous les} \ill{}, \ill{} est que pour \ill{} \ill{} \ill{}
coefficient $A = \mathbb{Z}/\ell\mathbb{Z}$, \add{p. ex. $A =
\mathbb{Z}/\ell^2\mathbb{Z}$} qui \ill{} \ill{} \ill{} \ill{} \ill{} des
\uncertain{limites} \ill{} \ill{}. [Je ne \uncertain{saurais} pas \ill{}
\ill{} \add{\uncertain{tel}} : dès que \ill{} n'est pas
$\mathbb{Z}/\ell\mathbb{Z}$, on ne sait pas \ill{} des résultats plus précis
que « le nb des pts fixes mod $\ell$ » — \uncertain{ainsi} pour le cas
\ill{} il suffirait \ill{} \ill{} de \ill{} \ill{} \uncertain{cohom.} des
faisceaux cst !].

On \uncertain{a} alors les $H^i(X, F)$ ne sont plus \uncertain{libres} sur
$A$, et les $\operatorname{Tr} f^{(i)}$ ne sont pas définis.

Une méthode, \uncertain{rédigée} dans \ill{} \ill{} (\uncertain{exposé
Bourbaki}) \uncertain{de tourner} la difficulté, est de \ill{} \ill{}
\ill{} \ill{} \uncertain{projection} \ill{} \ill{} \ill{} \ill{}
\ill{} \ill{} des faisceaux \ill{} \ill{} comme
$\mathbb{Z}/\ell^\nu\mathbb{Z}$ (\ill{} \ill{} \ill{}) en
\uncertain{passage} \ill{} \ill{} \ill{} \ill{} de coefficients tels
\ill{} $\mathbb{Z}_\ell$, \struck{\ill{}} qui \uncertain{sont} \ill{} des
anneaux \uncertain{intègres}. Cette \ill{} \ill{} \ill{} fonction
$\mathbb{Q}_\ell$ de car. $0$ !) — \ill{}

\page{127}
\uncertain{sont que en termes d'ensembles finis} \ill{} les \uncertain{modules}
de type fini sur le dit $A$ \ill{} \uncertain{définis} — et la formule de
Lefschetz \add{(\uncertain{formules de calcul de} \ill{})} \ill{} de
\uncertain{la} \ill{} \ill{} \ill{} \ill{} \uncertain{sous} \ill{}
\uncertain{exacte} : Cependant, \ill{} \ill{} \ill{} que \ill{} formule
\ill{} \uncertain{par} la \uncertain{conjugaison} \ill{} \ill{} \ill{} la
\uncertain{limite}, \ill{} \ill{} formules de Lefschetz par les \ill{}
\uncertain{écrites} par diagrammes \add{des} \ill{} \ill{} de coeff.
$\mathbb{Z}/\ell^\nu\mathbb{Z}$.

Comme l'a \uncertain{souligné} \ill{}, pour \uncertain{écrire} une \ill{}
\ill{} (pour un $A = \mathbb{Z}/N\mathbb{Z}$ \uncertain{donné}) il y a lieu
de remplacer les $H^i(X, A_X)$ (disons) par le
\uncertain{système} \add{des} $\mathbf{R}\Gamma_X(A_X)$ \add{qui lui
\uncertain{donne} naissance}, et les systèmes des $f^{(i)}$ par
l'endomorphisme $\mathbf{R}\Gamma_X(f)$ \ill{} \ill{} \ill{} complexe
\uncertain{défini} par « \ill{} ».

Pour ceci, considérons le \uncertain{formalisme} général \uncertain{suivant}
: un complexe $L_\bullet$ sur un anneau \uncertain{commutatif} $A$

\page{129}
(pas \uncertain{nécessairement} \uncertain{noethérien}). Nous supposons
$L_\bullet \in \operatorname{ob} D^-(A)$, et disons que $L_\bullet$ est
« pseudo-cohérent » s'il est isomorphe dans $D^-(A)$ à un complexe qui est
\add{\uncertain{d'$A$-modules}, \uncertain{ce qui} équivaut \uncertain{au
fait que les} $H_i(L_\bullet)$ \uncertain{de type fini}} \ill{} de type fini
\uncertain{en} tout degré ; \add{(\ill{} \ill{} \ill{} \ill{})} « \ill{}
\uncertain{dimension finie} » s'il est \uncertain{isomorphe} \ill{}
\struck{\ill{}} à un complexe \textit{plat} qui \ill{} \ill{} \ill{}
$\otimes$. [\uncertain{Dans} les \uncertain{complexes} \ill{} \ill{}
\ill{} \ill{} \ill{} \ill{} \ill{} \ill{} \uncertain{Variante} des
\ill{} \ill{} \ill{} catégories \uncertain{triangulées}]. Les complexes
pseudo-cohérents de Tor-dimension finie sont \uncertain{ceux} qui sont
isomorphes \add{\uncertain{par ex.}} dans $D^b(A)$ à des complexes qui
sont \textit{projectifs de type fini} en chaque degré, \textit{et nuls en
presque tous degrés}.

Si $u$ est un \uncertain{endomorphisme} d'un tel $L_\bullet$, on peut lui
associer une \uncertain{trace}

\page{131}
\uncertain{dite}
\[
  \operatorname{Tr}_* u \in A
\]
de telle façon que \textit{α)} si $L_\bullet$ est \textit{projectif} de type
fini \uncertain{en chaque} degré, \uncertain{et} \uncertain{nul} \ill{}
\ill{}, alors
\[
  \operatorname{Tr}_* u = \sum (-1)^i \operatorname{Tr} u_i ,
\]
\ill{} \textit{β)} $\operatorname{Tr}_* u$ est « stable par
isomorphismes » dans $D(A)$.

D'ailleurs, on \uncertain{peut} \uncertain{être} aussi
\uncertain{explicite} par l'ensemble des conditions
\note{les deux lettres sont permutées sur la page par une flèche : le b)
écrit en premier est renvoyé après le a).}

b) « \uncertain{additivité} » pour une \uncertain{suite exacte}
\ill{} \ill{} \ill{} ;

a) si $L_\bullet$ est \uncertain{réduit au degré} $0$ \ill{}
\ill{}, alors $\operatorname{Tr}_* u = \operatorname{Tr} u$.

\struck{Linéarité} \textit{D'autres} \uncertain{propriétés} formelles,
\uncertain{en outre} \ill{} :
c) $\operatorname{Tr} uv = \operatorname{Tr} vu$,

d) $\operatorname{Tr}\bigl(u \overset{\mathbf{L}}{\otimes} v\bigr)
= (\operatorname{Tr} u)(\operatorname{Tr} v)$,

qui se \uncertain{ramènent} aux formules correspondantes
\uncertain{connues} par ailleurs.

L'\textit{unicité} d'une fonction $\operatorname{Tr}$ satisfaisant aux
conditions α) et β) est claire. L'existence \struck{\ill{}} se ramène à
prouver que l'expression α) est

\page{133}
\textit{stable} par \uncertain{isomorphismes}, — qui \uncertain{semble}
\add{(et l'additivité et \uncertain{par} l'additivité habituelle)}
\uncertain{des} \uncertain{pseudo-cohérents}) se \uncertain{vérifie}
\ill{} \ill{} \ill{} \ill{} \ill{} \ill{} \ill{} \uncertain{l'additivité}
\ill{} \ill{} \ill{} \ill{} \ill{} \uncertain{tensoriel} \ill{}
\uncertain{trivialement}.

Pour usage \uncertain{ultérieur}, on peut présenter aussi la
\uncertain{définition} de $\operatorname{Tr}_* u$ comme \uncertain{suit},
\ill{} \ill{} \ill{} \uncertain{termes} \ill{}. Si $L_\bullet$ est
ps.\ coh.\ et de Tor-dim.\ finie, alors
\[
  \check{L}_\bullet \simeq \mathbf{R}\mathcal{H}om(L_\bullet, A),
\]
\uncertain{posons} — \ill{} \ill{} \ill{} ;
$\check{L}_\bullet$ est \uncertain{aussi} \add{ps.} \ill{} \ill{} \ill{}
$\operatorname{ob} D^-(A)$, \add{\uncertain{pseudo-cohérent}
\uncertain{et de Tor-dim.\ finie}} \uncertain{isomorphisme}
\struck{\ill{}} \ill{}, et \uncertain{le morphisme canonique}
\[
  L_\bullet \to \check{\check{L}}_\bullet
\]
est un \uncertain{isomorphisme}. D'ailleurs
\[
  \bigl(L_\bullet \overset{\mathbf{L}}{\otimes} L'_\bullet\bigr)^{\vee}
  \simeq \check{L}_\bullet \overset{\mathbf{L}}{\otimes} \check{L}'_\bullet .
\]
\uncertain{Enfin} \uncertain{pour} $M_\bullet \in \operatorname{ob} D^+(A)$,
\ill{} \ill{} \ill{} \uncertain{isomorphisme} \uncertain{canonique}
\[
  \mathbf{R}\mathcal{H}om^\bullet(L_\bullet, M_\bullet) \simeq
  \check{L}_\bullet \overset{\mathbf{L}}{\otimes} M_\bullet ,
\]
et si $M_\bullet$ \uncertain{satisfait} aux \uncertain{mêmes} conditions
\uncertain{que} $L_\bullet$, \uncertain{de} \uncertain{même}

\page{135}
\[
  \mathbf{R}\mathcal{H}om^\bullet(M_\bullet, L_\bullet) \simeq
  \check{M}_\bullet \overset{\mathbf{L}}{\otimes} L_\bullet ,
\]
et l'on voit que $\mathbf{R}\mathcal{H}om(M_\bullet, L_\bullet)$ et
$\mathbf{R}\mathcal{H}om(L_\bullet, M_\bullet)$ sont « \uncertain{duals} »
l'un de l'autre au sens de ${}^\vee$, en \uncertain{particulier}
\uncertain{ils sont} \textit{\ill{}} : \uncertain{chacun} \ill{} \ill{}
$A$. En \uncertain{particulier}, \ill{} \ill{} \uncertain{complexes}
\uncertain{résultats} \ill{} : si
\[
  u \in H^0\bigl(\mathbf{R}\mathcal{H}om^\bullet(L_\bullet, M_\bullet)\bigr)
  = \operatorname{Hom}_{D(A)}(L_\bullet, M_\bullet) ,
\]
\[
  v \in H^0\bigl(\mathbf{R}\mathcal{H}om(M_\bullet, L_\bullet)\bigr)
  = \operatorname{Hom}_{D(A)}(M_\bullet, L_\bullet) ,
\]
\ill{} \uncertain{accouplement} \uncertain{produit}
\[
  \langle u, v \rangle = H^0(A) = A .
\]

\struck{\uncertain{Que} \ill{}} \uncertain{Le} \uncertain{calcul}
\uncertain{donne} l'\uncertain{explicitation}
\[
  \langle u, v \rangle = \langle v, u \rangle
  = \sum (-1)^i \operatorname{Tr}(vu)^i = \sum (-1)^i \operatorname{Tr}(uv)^i
\]
\note{deux lettres sont surchargées dans la première somme ; on lit $(vu)$
puis $(uv)$ sans certitude.}
\uncertain{si} $L_\bullet$, $M_\bullet$ sont \struck{\ill{}}
\add{projectifs} de type \uncertain{fini} \add{\uncertain{en chaque
degré}} de \uncertain{type} \ill{} \uncertain{constants} (\uncertain{dans}
$K(A)$ !) \add{$u$, $v$ \uncertain{désignant} \ill{} \ill{} \ill{}}
\struck{\uncertain{qui} \ill{} \ill{} \ill{}} \ill{} \ill{}
\uncertain{chaque} \uncertain{degré}, \ill{} \ill{} \ill{}
$L_\bullet = M_\bullet$. \ill{} \uncertain{pour}
$u \in \operatorname{Hom}_{D(A)}(L_\bullet, L_\bullet)$ :
\[
  \sum (-1)^i \operatorname{Tr} u^i = \langle u, \mathrm{id}_{L_\bullet}
  \rangle = \langle \mathrm{id}_{L_\bullet}, u \rangle ,
\]
ce qui fournit bien une expression de $\sum (-1)^i \operatorname{Tr} u^i$
en termes \uncertain{intrinsèques}.
\marginal{b) \uncertain{Faire intervenir} \ill{} \ill{} \uncertain{des}
\ill{} soit plus \uncertain{général} \ill{} \ill{} …}

Bien entendu, \uncertain{pour} \uncertain{pouvoir} \uncertain{appliquer}
la \uncertain{notion} $\operatorname{Tr}_* u$ dans le
\struck{\uncertain{contexte}} \uncertain{cadre} de la formule de
\uncertain{Lefschetz}, il \uncertain{faut} \uncertain{savoir} que
$\mathbf{R}\Gamma_X(A_X)$ \uncertain{est} \add{\uncertain{ps.\ coh.\ et}}
de Tor-dimension \uncertain{finie}. \ill{} \ill{} \ill{} \ill{} \ill{}
\uncertain{les} \uncertain{standard}

\page{137}
\marginal{\textit{Remarques finales} : \uncertain{1°} Le formalisme des
$\operatorname{Tr}^*$ (et des « traces \uncertain{partielles} »)
\uncertain{marche} \uncertain{sur} un \uncertain{espace annelé}, avec
hypothèse pseudo-cohérente \ill{}}

\section{(31) Complexes de Tor-dimension finie sur un préschéma}

$X$ préschéma, $L_\bullet \in \operatorname{ob} D^-(X)$. On dit que
\struck{$X$} $L_\bullet$ est de Tor-dimension finie si pour tout
\struck{faisceau} \add{Module} $F$ sur $X$, on a
\[
  \mathcal{T}or_i(F, L_\bullet) = 0
\]
\add{(\uncertain{complexes} \ill{} \ill{}, par \uncertain{lui-même} !)}
pour \struck{\ill{}} $i$ grand. De façon précise, on dit que $L_\bullet$
est de \struck{\ill{}} Tor-dimension bornée par $[N, M]$,
\struck{\ill{}} $N, M \in \mathbb{Z}$, si
\[
  \mathcal{T}or_i(F, L_\bullet) = 0 \quad \text{pour } i \notin [N, M] .
\]
\struck{\uncertain{Cela posé},}

On voit que ceci équivaut à dire que pour tout pt géom.\ $\bar{x}$ de $X$,
le \uncertain{complexe} $L_{\bar{x}}$ est de Tor-dim.\ bornée par
$[N, M]$.

\struck{\ill{}}
\note{suit un bloc d'une dizaine de lignes, encadré et barré de traits
obliques, qui commençait par « Cela signifie également » et énonçait sous
a) une condition sur $L'_\bullet$ ; il n'est pas lu.}

Cela signifie aussi que si $L'_\bullet \in \operatorname{ob} D^b$ a sa
\struck{\ill{}} cohomologie bornée par $[N', M']$, alors
$L_\bullet \overset{\mathbf{L}}{\otimes} L'_\bullet$ a sa cohomologie bornée
par

\page{139}
\struck{\ill{}} $[N + N', M + M']$. Sous cette forme, il est évident que si
$L_\bullet$ est de Tor-dim.\ \uncertain{majorée} par $[N, M]$, et
$L'_\bullet$ de Tor-dim.\ majorée par $[N', M']$, alors
$L_\bullet \overset{\mathbf{L}}{\otimes} L'_\bullet$ est de Tor-dimension
majorée par $[N + N', M + M']$.
\add{On voit de même que si $f \colon X \to Y$, et $L_\bullet$ sur $Y$ est
de Tor-dim.\ majorée par $[N, M]$, alors $\mathbf{L}f^*(L_\bullet)$ itou
…}

\textit{Remarques.} a) Ces notions sont valables pour un faisceau
d'anneaux \textit{quelconque} sur $X$.

b) On voit facilement que $L_\bullet$ de Tor-dim.\ majorée par $[N, N']$
$\Longleftrightarrow$ \struck{\ill{}} $\exists\, L'_\bullet \to L_\bullet$
quasi-isom., avec $L'_\bullet$ \textit{plat} en chaque degré, et
$L'_i = 0$ si $i \notin [N, N']$. D'ailleurs, \ill{}, \uncertain{et}
$L_\bullet$ est à \uncertain{cohomologie} bornée par $[N, N']$ ssi
$L_\bullet$ est \struck{\ill{}} \uncertain{isomorphe} dans $D(X)$ à
$L'_\bullet$, \struck{\ill{}} \uncertain{avec} $L'_i = 0$ si
$i \notin [N, N']$ ($L'_\bullet$ \uncertain{non} plat !). Les
\uncertain{assertions} de \uncertain{stabilité} plus haut en résultent
formellement.

\textit{Proposition} \add{(Ici, il faut que l'Anneau de base soit de
torsion)}. Soit $f \colon X \to Y$ \uncertain{compactifiable} à fibres de
dim.\ $\leq d$. Si $F_\bullet \in \operatorname{ob} D^b(X)$ est de
Tor-dim.\ majorée par $[N, M]$, alors
\note{l'énoncé de la proposition se poursuit au-delà de ce lot.}

\end{document}
